% Source: https://artofproblemsolving.com/community/c3h2317020_number_theory
% Stable identity: aops-topic-2317020-q01
\textbf{(a)} Yes. Check: $7\cdot3=21=5\cdot4+1$ and $5=4+1$, so $7\cdot3\equiv5\pmod4$. The student's step is valid because $\gcd(7,4)=1$, so $7$ has an inverse modulo $4$: $7\cdot3=21\equiv1\pmod4$, so $7^{-1}\equiv3$.

Another way: reduce first. $7\equiv3$ and $5\equiv1\pmod4$, so $3x\equiv1\pmod4$. Since $3\cdot3=9\equiv1\pmod4$, multiplying by $3$ gives $x\equiv3\pmod4$. The solutions are $\boxed{x\equiv3\pmod4}$.

\textbf{(b)} If $\gcd(k,m)=1$, then $ka\equiv kb\pmod m$ implies $a\equiv b\pmod m$: indeed $m\mid k(a-b)$ and $m$ has no common factor with $k$, so $m\mid a-b$ (equivalently, multiply by the inverse of $k$). In general one may only conclude
\[a\equiv b\pmod{\dfrac{m}{\gcd(k,m)}}.\]
Example: $2\cdot1\equiv2\cdot3\pmod4$ (both sides are $\equiv2$), but $1\not\equiv3\pmod4$; the correct conclusion is $1\equiv3\pmod2$.
